\section{Introduction}
\label{sec:intro}

% TODO: open with clinical motivation paragraph (low back pain
% prevalence, MRI grading workload, inter-rater variability).

% TODO: state the gap — most spine grading models have rigid label
% spaces, so a model trained on RSNA's 5 stenosis labels cannot answer
% questions about Pfirrmann grading or Modic changes from SPIDER.

% TODO: state our contributions (3 bullets).

\textbf{Contributions.}
\begin{itemize}
  \item A 3D CBAM-augmented backbone that recovers minority Severe-class
        detection on imbalanced spinal stenosis grading
        (Foraminal Severe F1 from $0.000 \to 0.18/0.20$).
  \item A hybrid integration of frozen BiomedCLIP image and text
        encoders with the 3D backbone, enabling zero-shot extension
        to unseen disease labels via cosine-similarity classification.
  \item A cross-dataset zero-shot evaluation (RSNA $\to$ SPIDER) showing
        \emph{selective transfer}: the hybrid wins on disc-related
        labels where 3D context matters, while a 2D-only baseline wins
        on vertebra/endplate appearance — characterising when the
        added complexity is justified.
\end{itemize}

% Motivation for the design choice (vs MedGemma / RadFM) is summarised
% in MOTIVATION_AND_DESIGN.md (repo); short version belongs in
% Section~\ref{sec:related} of this paper.
